\chapter{Conclusões e Trabalhos Futuros}
\label{cap:conclusoes}

% TODO(P1): Sintetizar as conclusões finais do estudo, cobrindo:
%
% 1. Síntese dos Resultados:
%    - Confirmação empírica do ganho assintótico da Programação Dinâmica em relação a abordagens ingênuas.
%    - Papel essencial do reconhecimento prévio da subestrutura ótima e sobreposição de subproblemas.
%
% 2. Onde a Programação Dinâmica Falha ou é Limitada:
%    - Ausência de Subestrutura Ótima: problema do caminho simples mais longo em grafos (subproblemas não são independentes).
%    - Explosão do Espaço de Estados: Caixeiro Viajante (TSP) com o algoritmo de Held-Karp atingindo O(n^2 \cdot 2^n).
%    - Natureza Pseudo-Polinomial: Problema da Mochila 0/1 com complexidade O(n \cdot W), onde W é o valor numérico.
%    - Gargalo de Memória: problemas com matrizes multidimensionais onde o espaço Theta(n^2) ou Theta(n^3) esgota a RAM.
%
% 3. Trabalhos Futuros:
%    - Paralelização de preenchimento de tabelas em GPUs/multithreading.
%    - Heurísticas aproximativas para problemas NP-difíceis onde a PD clássica é inviável.

% [Texto a cargo de P1 - ver docs/tarefas/P1_relatorio_secoes_1_2_6.md]
